\documentclass[11pt,a4paper]{article}
\usepackage[T1]{fontenc}
\usepackage[utf8]{inputenc}
\usepackage{lmodern}
\usepackage{amsmath,amssymb,amsthm}
\usepackage{booktabs}
\usepackage{xcolor}
\usepackage[margin=27mm]{geometry}
\usepackage[colorlinks=true,urlcolor=blue,citecolor=blue,linkcolor=blue]{hyperref}
\newtheorem{theorem}{Theorem}
\newtheorem{proposition}[theorem]{Proposition}
\newtheorem{lemma}[theorem]{Lemma}
\newtheorem{corollary}[theorem]{Corollary}
\theoremstyle{definition}
\newtheorem{definition}[theorem]{Definition}
\theoremstyle{plain}
\DeclareMathOperator{\im}{im}
\DeclareMathOperator{\rank}{rank}
\newcommand{\F}{\mathbb{F}_2}
\newcommand{\Z}{\mathbb{Z}}
\newcommand{\T}{\mathsf{T}}
\newcommand{\tQ}{\tilde Q}
\newcommand{\hp}{\hat p}
\newcommand{\htau}{\hat\tau}
\newcommand{\code}[1]{[\![#1]\!]}
\setlength{\emergencystretch}{2em}
\setlength{\parindent}{0pt}
\setlength{\parskip}{0.3em plus 0.1em minus 0.1em}
\title{Covers of bivariate bicycle codes:\\
partial results on the distance conjecture for power-of-two covers}
\author{Fable \& Astra\\[0.35em]
\normalsize Research led by Jean-Philippe Bruneton\\[-0.1em]
\small Universit\'e Paris Cit\'e, France}
\date{19 November 2026}

\begin{document}
\maketitle
\begin{abstract}
Bivariate bicycle codes are quantum error-correcting codes whose qubits and
checks sit on a discrete torus. Enlarging the torus by an integer factor $h$
produces a larger code whose Tanner graph is an $h$-fold cover of the
original one, and Symons, Rajput and Browne conjectured that such a cover
never shortens the minimum distance and never multiplies it by more than $h$:
$d\le d_h\le hd$. They proved the upper bound for odd $h$ and the lower bound
for odd $h$ when the number of logical qubits does not change. We treat
covers whose degree is a power of two, where their arguments do not apply. A
double cover gives a short exact sequence of chain complexes over $\F$, and
its long exact sequence in homology yields three statements. The number of
logical qubits satisfies $k\le k_h\le hk$, and for a double cover the kernel
of the transfer on logical operators has dimension exactly $k-k_2/2$. The
distance satisfies $d_h\ge\min(d,s)$, where $s$ is the least weight of a
nonzero stabiliser of the base code; this proves the conjectured lower bound
for every power-of-two cover of a base code with $s\ge d$, which includes
805 of the 836 inequivalent base codes with $k>0$ on tori of at most $27$ sites, and the
$\code{72,12,6}$ code. Finally $d_2\le2d$ whenever $k_2=2k$. An example
$\code{42,4,6}\to\code{84,4,10}$ shows that the transfer alone cannot prove
the upper bound in general. Exact scans of about $2\times10^5$ covers of small
codes find no counterexample to the conjecture.
\end{abstract}

\section{Introduction}
A bivariate bicycle code is a quantum low-density parity-check code built on
a discrete torus $\Z_l\times\Z_m$. Each site of the torus carries two qubits,
one $X$-type check and one $Z$-type check, and every check acts on six qubits
selected by six fixed translations. Bravyi and coauthors showed that such
codes combine a high error threshold with a much smaller qubit overhead than
the surface code~\cite{Bravyi2024}; their best-known member is the
$\code{144,12,12}$ ``gross'' code.

Enlarging the torus in a compatible way produces a new code of the same family
that projects onto the original one: its Tanner graph is a covering graph of
the original Tanner graph. Symons, Rajput and Browne~\cite{Symons2025} turned
this observation into a systematic construction of sequences of bivariate
bicycle codes, and the gross code itself is a double cover of the
$\code{72,12,6}$ code, with exactly twice its distance. The central question
is how covering changes the minimum distance, the least weight of a
nontrivial logical operator. Topological intuition suggests the answer: an
$h$-sheeted cover neither shortens a nontrivial cycle nor lengthens it by more
than a factor $h$. Symons, Rajput and Browne proved parts of this, for odd $h$,
and conjectured the rest:

\emph{Conjecture~\cite{Symons2025}.} Every $h$-cover of a bivariate bicycle
code satisfies $d\le d_h\le hd$.

Their arguments rest on the fact that, for odd $h$, the composition of the
projection onto the base with the transfer back to the cover is multiplication
by $h$, which is invertible over $\F$. For even $h$ this composition vanishes,
and both bounds were left open; the lower bound was also open for odd $h$ when
the cover has more logical qubits than the base.

This note does not settle the conjecture, but it proves three new statements
for covers of degree a power of two. For a double cover, the vanishing of
projection after transfer turns into an exact sequence: the cover is an
extension of the base by itself. The resulting long exact sequence in
homology gives the bounds $k\le k_2\le2k$ on the number of logical qubits and
the exact dimension $k-k_2/2$ of the space of base logical operators killed
by the transfer (Theorem~\ref{thm:k}). It also gives a lower bound on the
distance, $d_2\ge\min(d,s)$ with $s$ the least weight of a nonzero
stabiliser, which survives iteration along towers of double covers and proves
$d\le d_h$ for every power-of-two cover when $s\ge d$
(Theorem~\ref{thm:lower}). Finally $d_2\le2d$ when $k_2=2k$
(Theorem~\ref{thm:upper}). The example of Section~\ref{sec:example} shows why
the upper bound needs more than the transfer, and the scans of
Section~\ref{sec:scans} test the conjecture on small codes.

\section{Related work}\label{sec:related}
Bivariate bicycle codes were introduced by Bravyi et al.~\cite{Bravyi2024},
who also established the two structural facts used below: the check matrices
$H_X$ and $H_Z$ have equal rank, and the $X$- and $Z$-distances coincide.

The conjecture and the covering framework are due to Symons, Rajput and
Browne~\cite{Symons2025}. They define $h$-covers by an algebraic condition on
the defining polynomials, show that projection and transfer are chain maps,
and prove $k_h\ge k$ for all $h$, $d_h\le hd$ for odd $h$, and $d\le d_h$ for
odd $h$ when $k_h=k$. Our Theorems~\ref{thm:k}--\ref{thm:upper} are the even
counterpart of these statements for degrees that are powers of two; the upper
bound $k_h\le hk$ of Theorem~\ref{thm:k} has no analogue in their paper.

Hirasaki and Lee~\cite{Hirasaki2026} extend the odd-index statements to lifts
of lifted product codes: for odd lift index $t$ they obtain $k'\ge k$,
$d'\le td$, and $d'\ge d$ when $k'=k$, and they note that the transfer need
not be injective for even $t$ and that no analogous relation follows directly
in that case. Our Lemma~\ref{lem:exact} is the tool that replaces injectivity
for index two. Sabo, Can and Marquis~\cite{Sabo2026} study bivariate bicycle
codes over finite fields through the orbit decomposition of the defining
ring, give an independent algebraic proof of dimension lower bounds for cover
codes, an exact dimension formula, and distance certificates; their methods
are algebraic and do not use the exact sequence below.

The sequence of Lemma~\ref{lem:exact} is the chain-level form of the transfer
exact sequence of a two-sheeted covering space with $\F$ coefficients, a
standard tool of algebraic topology. We did not find its application to
bivariate bicycle codes in the documents examined (Section~\ref{sec:scope}).

\section{Setting}\label{sec:setting}
Let $R=\F[x,y]/(x^l-1,\,y^m-1)$, identified with the group algebra of
$\Z_l\times\Z_m$, and let $A,B\in R$ be sums of three distinct monomials.
Each element of $R$ acts on $\F^{N}$, $N=lm$, by multiplication, and we write
$A$ and $B$ also for these $N\times N$ matrices. The code $Q(A,B,l,m)$ is the
Calderbank--Shor--Steane code with check matrices $H_X=(A\ \ B)$ and
$H_Z=(B^\T\ \ A^\T)$. It has $n=2N$ qubits, $k$ logical qubits and distance
$d$; distances are defined only when $k>0$, and we read the conjecture for
base codes with $k>0$.

The code is the chain complex
\[
Q_2=\F^{N}\xrightarrow{\ \partial_2=H_Z^\T\ }Q_1=\F^{2N}
\xrightarrow{\ \partial_1=H_X\ }Q_0=\F^{N}.
\]
A \emph{cycle} is an element of $\ker\partial_1$, that is, a $Z$-type
operator commuting with all $X$ checks, and a \emph{boundary} is an element of
$\im\partial_2$, a $Z$-type stabiliser. Nontrivial $Z$-type logical operators
are cycles that are not boundaries, $k=\dim H_1(Q)$, and $d$ is the least
weight of a nontrivial cycle (the $X$- and $Z$-distances coincide). We write
$s$ for the least weight of a nonzero boundary. Since each $Z$ check has
weight six, $s\le6$. Equality of the ranks of $H_X$ and $H_Z$ gives
$k=2N-2\rank H_Z$, hence $\dim H_2(Q)=\dim\ker\partial_2=N-\rank H_Z=k/2$.

A code $\tQ=Q(\tilde A,\tilde B,ul,tm)$ is an \emph{$h$-cover} of $Q$, with
$h=ut$, when reducing the exponents of $x$ modulo $l$ and of $y$ modulo $m$
maps the three monomials of $\tilde A$ onto those of $A$ and those of
$\tilde B$ onto those of $B$~\cite{Symons2025}. Its parameters are
$\code{hn,k_h,d_h}$, and its $s$-parameter is $s_h$. The \emph{projection}
$p:\tQ_i\to Q_i$ adds the coordinates of each fibre of the reduction map, and
the \emph{transfer} $\tau=p^\T:Q_i\to\tQ_i$ copies each coordinate of the base
onto its whole fibre. Symons, Rajput and Browne show that $p$ and $\tau$ are
chain maps and that $p\circ\tau=h$~\cite{Symons2025}. We write $\hp$ and
$\htau$ for the induced maps in homology. Projection can only lower the
weight, $|p(v)|\le|v|$, and $|\tau(v)|=h|v|$. A \emph{double cover} is a
$2$-cover, $(u,t)=(2,1)$ or $(1,2)$; a tilde or an index $2$ refers to it.

\section{Results}\label{sec:results}
\subsection{The exact sequence of a double cover}

\begin{lemma}\label{lem:exact}
For a double cover, $0\to Q\xrightarrow{\ \tau\ }\tQ\xrightarrow{\ p\ }Q\to0$
is a short exact sequence of chain complexes. It induces the long exact
sequence
\[
0\to H_2(Q)\xrightarrow{\htau}H_2(\tQ)\xrightarrow{\hp}H_2(Q)
\xrightarrow{\delta}H_1(Q)\xrightarrow{\htau}H_1(\tQ)\xrightarrow{\hp}H_1(Q)
\xrightarrow{\delta'}H_0(Q)\to\cdots
\]
\end{lemma}

\begin{proof}
In each degree $\tau$ is injective, $p$ is surjective, and
$p\circ\tau=2=0$, so $\im\tau\subseteq\ker p$. The dimensions agree:
$\dim\ker p=2\dim Q_i-\dim Q_i=\dim\im\tau$. Hence $\ker p=\im\tau$, and the
long exact sequence is the one associated with any short exact sequence of
chain complexes (snake lemma); it starts with $0$ because the complexes vanish
in degree~$3$.
\end{proof}

\subsection{Logical qubits and the kernel of the transfer}

\begin{theorem}\label{thm:k}
For a double cover, $k\le k_2\le2k$ and
$\dim\ker\bigl(\htau:H_1(Q)\to H_1(\tQ)\bigr)=k-k_2/2$. For an $h$-cover with
$(u,t)=(2^a,2^b)$, $h=2^{a+b}$, one has $k\le k_h\le hk$.
\end{theorem}

\begin{proof}
By exactness, $\ker\htau=\im\delta$ in degree one. The initial segment
$0\to H_2(Q)\to H_2(\tQ)\to H_2(Q)$ is exact, and $\dim H_2(\tQ)=k_2/2$ by the
rank equality for the cover, so
$\dim\im\bigl(\hp:H_2(\tQ)\to H_2(Q)\bigr)=k_2/2-k/2$. Hence
$\dim\im\delta=\dim H_2(Q)-\dim\im\hp=k/2-(k_2/2-k/2)=k-k_2/2$. This
dimension lies between $0$ and $\dim H_2(Q)=k/2$, which gives
$k\le k_2\le 2k$.

A $(2^a,2^b)$-cover factors into $a+b$ successive double covers through the
tori $(2^il,2^jm)$; each intermediate code is the bivariate bicycle code
obtained by reducing the exponents of $\tilde A$ and $\tilde B$ modulo the
intermediate sizes. Its three monomials stay distinct, since they reduce to
the distinct monomials of the base, so each step is a double cover in the
above sense and the bounds multiply.
\end{proof}

The argument does not give $k_h\le hk$ for odd $h$.

\subsection{A lower bound on the distance}

\begin{theorem}\label{thm:lower}
For a double cover, $d_2\ge\min(d,s)$ and $s_2\ge\min(s,2d)$. Consequently
$d_h\ge\min(d,s)$ for every $(2^a,2^b)$-cover. In particular the conjectured
bound $d\le d_h$ holds for every cover of degree a power of two of a base code
with $s\ge d$.
\end{theorem}

\begin{proof}
Let $\tilde L$ be a cycle of the double cover that is not a boundary. Its
projection $p(\tilde L)$ is a cycle of the base, and there are three cases.
If $p(\tilde L)$ is not a boundary, $|\tilde L|\ge|p(\tilde L)|\ge d$. If
$p(\tilde L)$ is a nonzero boundary, $|\tilde L|\ge|p(\tilde L)|\ge s$. If
$p(\tilde L)=0$, Lemma~\ref{lem:exact} gives $\tilde L=\tau(L)$ for some
$L\in Q_1$. Since $\tau$ is an injective chain map,
$\tau(\partial_1L)=\tilde\partial_1\tilde L=0$ forces $\partial_1L=0$, and
$L$ is not a boundary, because $L=\partial_2z$ would give
$\tilde L=\tilde\partial_2\tau(z)$. Hence $|\tilde L|=2|L|\ge2d$. In all cases
$|\tilde L|\ge\min(d,s)$.

Let $\tilde b$ be a nonzero boundary of the cover. Its projection is a
boundary of the base. If it is nonzero, $|\tilde b|\ge s$. Otherwise
$\tilde b=\tau(c)$ with $c\ne0$, and $c$ is a cycle as above, either a
nonzero boundary or a nontrivial cycle, so $|c|\ge\min(d,s)$ and
$|\tilde b|\ge2\min(d,s)$. Thus $s_2\ge\min(s,2d,2s)=\min(s,2d)$.

For a tower of double covers with parameters $(d^{(j)},s^{(j)})$, starting
from $(d^{(0)},s^{(0)})=(d,s)$, put $\mu=\min(d,s)$. If $d^{(j)}\ge\mu$ and
$s^{(j)}\ge\min(s,2d)$, then
$d^{(j+1)}\ge\min(d^{(j)},s^{(j)})\ge\min(\mu,s,2d)=\mu$ and
$s^{(j+1)}\ge\min(s^{(j)},2d^{(j)})\ge\min(s,2d,2\mu)=\min(s,2d)$. Each
intermediate code has $k^{(j)}\ge k>0$ by Theorem~\ref{thm:k}, so its
distance is defined, and induction concludes.
\end{proof}

The proof locates any counterexample to the lower bound for a double cover:
it must be a nontrivial cycle of the cover of weight less than $d$ whose
projection is a nonzero stabiliser of the base. This requires $s<d$, and since
$s\le6$ and distances of these codes are even in all cases we computed, in
practice $d\ge8$.

The condition $s\ge d$ was computed exactly for all base codes with $k>0$ and
$lm\le27$, up to the symmetries listed in Section~\ref{sec:scans}. Of the
$836$ codes obtained, the $805$ codes with $d\le6$ all satisfy $s\ge d$, and the
remaining $31$ have $d=8$ and $s=6$. For the first $805$, the lower bound of the
conjecture is therefore a theorem for all covers of degree a power of two. The
$\code{72,12,6}$ code also satisfies $s=d=6$, so every power-of-two cover of it,
the gross code included, has distance at least six. The codes
$\code{90,8,10}$ and $\code{108,8,10}$ of~\cite{Bravyi2024} have $s\le6<d$ and
are not covered.

\subsection{An upper bound when the logical dimension doubles}

\begin{theorem}\label{thm:upper}
For a double cover with $k_2=2k$, the transfer $\htau$ is injective on
$H_1(Q)$ and $d_2\le2d$. More generally, for every double cover
$d_2\le2d_\tau$, where $d_\tau$ is the least weight of a cycle of the base
whose class lies outside the subspace $\im\delta=\ker\htau$, of dimension
$k-k_2/2\le k/2$.
\end{theorem}

\begin{proof}
If $L$ is a cycle of the base whose class is not in $\ker\htau$, then
$\tau(L)$ is a nontrivial cycle of the cover of weight $2|L|$. When
$k_2=2k$, Theorem~\ref{thm:k} gives $\ker\htau=0$, and $L$ can be taken to be
a nontrivial cycle of weight $d$.
\end{proof}

\subsection{An example where the transfer is not enough}\label{sec:example}
The transfer cannot prove the upper bound for even $h$ on its own. Take the
base $l=1$, $m=21$, $A=1+y+y^5$, $B=1+y+y^8$, with parameters
$\code{42,4,6}$, and its double cover $\tilde l=2$, $\tilde A=1+y+xy^5$,
$\tilde B=1+xy+xy^8$, with parameters $\code{84,4,10}$; both distances are
exact. Here $k_2=k$, the kernel of the transfer has dimension $k-k_2/2=2$, and
it contains the classes of all logical operators of weight six of the base:
the lightest base cycle with nontrivial transfer has weight $d_\tau=8$, so
Theorem~\ref{thm:upper} gives only $d_2\le16$. The conjectured bound
$d_2\le12$ nevertheless holds, since $d_2=10$. A logical operator of weight
ten of the cover is not a transfer, because it would be the transfer of a
nontrivial base cycle of weight five, below $d$. A proof of the upper bound
for even $h$ must therefore control the weight of lifts of the classes in
$\im\delta$, which the exact sequence does not do.

\subsection{Exact search for counterexamples}\label{sec:scans}
Three families of scans, all with exact distances, found no violation of
$d\le d_h\le hd$ or of $k_h\ge k$. Their outcomes are recorded in the data
file of the entry (Section~\ref{sec:verification}).

The first scan enumerates all pairs $(A,B)$ of trinomials that contain the
monomial $1$ (a normalisation without loss of generality, by translation) on
all tori with $l\le m$, and all lifts of their monomials in both directions.
It covers $38\,560$ double covers of bases with $lm\le12$, $21\,546$ triple
covers with $lm\le9$, and $39\,168$ fourfold covers of types $(4,1)$, $(1,4)$
and $(2,2)$ with $lm\le8$, together with $20\,992$ covers of type $(2,2)$ of
bases on the torus $(1,9)$; the distance of every cover is computed exactly.

The second scan works up to symmetry (translations, automorphisms
$x\mapsto x^a$ and $y\mapsto y^b$, exchange of $A$ and $B$) and over all
orientations $(l,m)$. It covers $3\,104$ double covers for $lm\le15$,
$17\,088$ for $16\le lm\le24$, $6\,240$ for $25\le lm\le27$, and $52\,974$
triple covers for $lm\le21$, of which $2\,015$ have $k_h>k$, the odd case left
open by~\cite{Symons2025}. Base codes reach distance $8$ ($n=48$ and $n=54$).
For each cover the absence of a nontrivial logical operator of weight below
$d$ is established by exhaustive enumeration, and the upper bound by the
transfer of an operator of weight $d$ when possible, by an exact computation
of $d_h$ otherwise.

The third scan treats the $96$ double covers of the codes $\code{72,12,6}$,
$\code{90,8,10}$ and $\code{108,8,10}$ of~\cite{Bravyi2024}, sixteen lifts per
direction. None has a nontrivial logical operator of weight below $d$
(exhaustive enumeration up to weight $d-1$), which tests the lower bound in the
regime $s<d$ that Theorem~\ref{thm:lower} does not reach. The upper bound
follows from the transfer for $94$ of them; for the remaining two, covers
$\code{216,8,d_h}$ of $\code{108,8,10}$, explicit logical operators of weights
$10$ and $12$ give $d_h\le12$, and the public checker shows that both have
$d_h=10$ exactly. The $486$ triple covers of these three codes all have
$k_h=k$ and fall under the theorems of~\cite{Symons2025}.

\section{Verification}\label{sec:verification}
Theorems~\ref{thm:k}, \ref{thm:lower} and~\ref{thm:upper} are proved above.
The public
\href{https://github.com/jpbruneton/open_quantum_problems/blob/main/research/entries/bivariate-bicycle-coverings/checks/check_bb_covers.py}{Python checker},
with its module of
\href{https://github.com/jpbruneton/open_quantum_problems/blob/main/research/entries/bivariate-bicycle-coverings/checks/bb_tools.py}{$\F$ tools},
tests them by exact computation, writes a
\href{https://github.com/jpbruneton/open_quantum_problems/blob/main/research/entries/bivariate-bicycle-coverings/data/verification.json}{verification record},
and reads the
\href{https://github.com/jpbruneton/open_quantum_problems/blob/main/research/entries/bivariate-bicycle-coverings/data/witnesses_108_covers.json}{certificate}
of the two explicit operators above. The module computes distances in three
independent ways: complete enumeration of the kernel, meet in the middle on
pairs (syndrome, logical value), and an anchored enumeration that exploits
translation invariance and the two-block structure.

The checker, which runs in about one minute, first validates the distance
routines on six known codes, $\code{18,8,2}$, $\code{36,8,4}$,
$\code{54,8,6}$, $\code{72,12,6}$, $\code{90,8,10}$ and $\code{108,8,10}$, by
comparing the available methods and checking $d_X=d_Z$ where the size allows.
It verifies that $p$ and $\tau$ are chain maps with $\ker p=\im\tau$ on the
cover $\code{36,8,4}\to\code{18,8,2}$. It then runs through the $672$ double
covers of the $42$ base codes with $lm\le9$, up to symmetry, and checks on each
one Theorems~\ref{thm:k}, \ref{thm:lower} and~\ref{thm:upper}: the bounds
$k\le k_2\le2k$, the dimension of the kernel of the transfer computed
directly, $d_2\ge\min(d,s)$, and $d_2\le2d$ with injective transfer in the
$50$ cases where $k_2=2k$; one cover in seventeen also has its number of
logical qubits recomputed by an independent dense elimination. It checks
$k_4\le4k$ and $d\le d_4\le4d$ on $1\,280$ fourfold covers, recomputes the
example $\code{42,4,6}\to\code{84,4,10}$ with $d_\tau=8$, shows that a double
cover of $\code{90,8,10}$ has no nontrivial logical operator of weight at most
nine, validates the two certificate operators and the bound $d_h\ge10$ for the
covers of $\code{108,8,10}$, and computes $s=6$ for $\code{72,12,6}$. The large
scans of Section~\ref{sec:scans} took several hours and are summarised in the
\href{https://github.com/jpbruneton/open_quantum_problems/blob/main/research/entries/bivariate-bicycle-coverings/data/scan_summary.json}{scan record}
of the entry; they are not rerun by the checker.

\section{Scope and limitations}\label{sec:scope}
\emph{Analytic proof.} Theorems~\ref{thm:k}, \ref{thm:lower}
and~\ref{thm:upper} hold for covers of degree a power of two. Their proofs use
only the equality of the ranks of $H_X$ and $H_Z$, valid for every bivariate
bicycle code, and the chain maps $p$ and $\tau$.

\emph{Computational verification.} The checker confirms these statements on
the small codes listed above without replacing the proofs. The scans are an
exact but finite search: they concern bases with $n\le54$ when systematic,
plus three named codes.

\emph{What remains open.} The upper bound $d_h\le hd$ for even $h$ is proved
only when $k_2=2k$, and the example of Section~\ref{sec:example} shows that the
transfer cannot give it in general. The lower bound $d\le d_h$ for degrees a
power of two is proved only when $s\ge d$, which excludes every base code of
distance above six; for those the evidence is numerical and limited to $31$
small codes of distance eight and three codes of~\cite{Bravyi2024}. For odd
$h$ with $k_h>k$, and for even $h$ that are not powers of two, nothing is
proved here; the evidence reduces to the $2\,015$ triple covers with $k_h>k$ of
the second scan, all from base codes of distance at most six. The qudit
version of the conjecture was not examined.

\emph{Bibliographic priority.} The search, on 19 November 2026, combined web
queries on the conjecture with the list of works citing~\cite{Symons2025}
returned by the Semantic Scholar interface (eighteen entries); the closest
ones were read at the source. Two listed titles could not be located, and
works on lifts of general Calderbank--Shor--Steane codes were not surveyed.
Theorems~\ref{thm:k}--\ref{thm:upper} were not found in the documents read.
They rest on a classical exact sequence and may exist elsewhere; priority is
claimed only within the limits of this search.

\begingroup\small
\begin{thebibliography}{9}
\setlength{\itemsep}{2pt}
\bibitem{Symons2025}
B. C. B. Symons, A. Rajput and D. E. Browne,
\emph{Sequences of Bivariate Bicycle Codes from Covering Graphs},
\href{https://arxiv.org/abs/2511.13560}{arXiv:2511.13560} (v1 2025, v3 2026).

\bibitem{Bravyi2024}
S. Bravyi, A. W. Cross, J. M. Gambetta, D. Maslov, P. Rall and T. J. Yoder,
\emph{High-threshold and low-overhead fault-tolerant quantum memory},
Nature \textbf{627}, 778--782 (2024);
\href{https://doi.org/10.1038/s41586-024-07107-7}{doi:10.1038/s41586-024-07107-7};
\href{https://arxiv.org/abs/2308.07915}{arXiv:2308.07915}.

\bibitem{Hirasaki2026}
Y. Hirasaki and J. Y. Lee, \emph{Lifting Lifted Product Codes},
\href{https://arxiv.org/abs/2607.28621}{arXiv:2607.28621} (2026).

\bibitem{Sabo2026}
E. Sabo, M. B. Can and D. Marquis,
\emph{Spectral Theory of Semisimple Bivariate Bicycle Codes},
\href{https://arxiv.org/abs/2608.27565}{arXiv:2608.27565} (2026).
\end{thebibliography}
\endgroup
\end{document}
